\chapter{三个月刷题计划}

三个月计划分成十二周。每周有一个算法主题、一个数据结构主题、一个 C++ 主题和一组力扣题。建议总题量控制在 180 到 260 题之间。题量不是越多越好；如果每道题都只是看题解，做 500 题也可能没有形成能力。更合理的目标是：核心题反复做，变形题系统做，难题适量做。

第一周学习复杂度、暴力法、数组和字符串。目标是建立“先写暴力，再找重复工作”的习惯。C++ 重点是 \code{vector}、\code{string}、下标边界、引用传参、\code{sort} 和 \code{lower_bound}。推荐题包括两数之和、移动零、删除有序数组重复项、合并两个有序数组、最大子数组和、反转字符串和有效回文。

第二周学习双指针和滑动窗口。目标是能判断窗口什么时候扩张、什么时候收缩。滑动窗口只适合窗口状态可维护且边界移动具有单调性的题；如果数组有负数，很多基于窗口和的算法会失效。推荐题包括盛最多水的容器、三数之和、长度最小的子数组、无重复字符的最长子串、找到字符串中所有字母异位词和最小覆盖子串。

第三周学习哈希表、集合和前缀和。目标是理解“用空间换查询时间”。哈希表优化的是反复查历史信息；前缀和优化的是反复求区间信息。推荐题包括字母异位词分组、最长连续序列、和为 K 的子数组、连续数组、快乐数和四数相加 II。

第四周学习栈、队列、堆和单调结构。栈处理最近未匹配，队列处理层次扩展，堆处理动态极值，单调结构处理可淘汰候选。推荐题包括有效括号、最小栈、每日温度、柱状图最大矩形、滑动窗口最大值、数组中的第 K 个最大元素和前 K 个高频元素。

第五周学习二分查找。目标是把二分从“找数组里的值”扩展到“找答案边界”。二分成立的关键是判断函数单调，而不只是数组有序。推荐题包括搜索插入位置、搜索旋转排序数组、寻找峰值、查找第一个和最后一个位置、爱吃香蕉的珂珂和分割数组的最大值。

第六周学习链表和树。链表训练指针修改顺序，树训练层次、递归关系和状态传递。本书会讲递归思想，但 C++ 代码会尽量给出迭代版本或说明递归栈风险。推荐题包括反转链表、环形链表、删除倒数第 N 个节点、二叉树层序遍历、验证二叉搜索树和最近公共祖先。

第七周学习图、BFS、DFS 和拓扑排序。很多图题难点不在搜索，而在建模：格子是节点，相邻关系是边；课程依赖是有向边；单词转换是隐式图。推荐题包括岛屿数量、腐烂的橘子、被围绕的区域、课程表、课程表 II 和单词接龙。

第八周学习回溯。目标是掌握选择、推进、撤销选择和剪枝。推荐题包括全排列、子集、组合总和、括号生成、单词搜索、N 皇后和分割回文串。

第九周学习一维动态规划。目标是理解状态定义。不要先背转移方程，先问 \code{dp[i]} 表示什么、它覆盖哪些决策、是否足够推出未来。推荐题包括爬楼梯、打家劫舍、零钱兑换、最长递增子序列和单词拆分。

第十周学习二维动态规划和背包。目标是掌握序列对齐、路径计数、容量约束和滚动数组。推荐题包括不同路径、最小路径和、最长公共子序列、编辑距离、分割等和子集和目标和。

第十一周学习贪心、并查集和高级结构。贪心必须证明局部选择为什么推出全局最优；并查集用于动态连通性；字典树用于前缀；树状数组和线段树用于区间查询更新。推荐题包括跳跃游戏、加油站、合并区间、冗余连接、账户合并和实现 Trie。

第十二周做模拟面试和查漏补缺。每天做一次 45 分钟模拟：5 分钟澄清题意，10 分钟讲暴力和优化，25 分钟写代码，5 分钟测试和复杂度总结。复盘时必须记录卡点，不要只写“没想到”。

\begin{principlebox}{每周交付}
每周至少提交一份报告，包含完成题目、三道重点复盘题、每类题的模板、暴力到优化的推理链、C++ 容器坑点和下周回访题。刷题报告不是形式，它是把题解转化成长期能力的工具。
\end{principlebox}
